\section{Financial Analysis Foundations}
\label{sec:financial-analysis-foundations}

Stock analysis combines several forms of evidence that answer different questions. Technical analysis examines how price and trading activity are evolving, fundamental analysis assesses the business and its valuation, and financial-news analysis identifies recent events that may change expectations. None of these perspectives is sufficient by itself. Stockrium therefore treats them as complementary inputs whose relevance varies with the investment horizon and whose outputs are synthesized by the multi-agent workflow described in Chapter~4. The same separation also makes configurations inspectable: users can see which evidence sources were enabled and compare how different supported selections and weights affect the resulting assessment.

\subsection{Technical Analysis}
\label{subsec:technical-analysis-foundations}

\paragraph{Concept and purpose.}

Technical analysis studies historical price and trading-volume behavior to describe how market participants are currently valuing and trading a security \cite{murphy1999technical}. Its principal concepts are \textit{trend}, the prevailing direction of price; \textit{momentum}, the strength and rate of a price movement; \textit{volatility}, the degree of price dispersion; \textit{volume}, the level of market participation; and \textit{support}, \textit{resistance}, and possible \textit{turning points}, which provide context for where buying or selling pressure may change. Charts represent the raw observations, while indicators transform those observations into comparable measures of trend, momentum, or volatility.

The purpose of technical analysis is to characterize the current market state, recognize changes in that state, and support the timing and risk management of a decision. It helps an analyst compare current behavior with recent history, identify potential entry or exit conditions, and judge whether a fundamentally or news-supported view is also reflected in price and volume. It does not estimate intrinsic value, explain the underlying business, or guarantee that a familiar pattern will be followed by the same outcome. Technical evidence is therefore most useful as conditional and time-sensitive decision support rather than as a deterministic prediction.

\paragraph{OHLCV market data.}

A price series is divided into fixed intervals, or \textit{bars}. Let $p_{t,j}$ and $q_{t,j}$ denote the price and traded quantity of transaction $j$ during interval $t$. The five basic values of a bar are

\begin{equation}
    O_t=p_{t,1}, \qquad
    H_t=\max_j p_{t,j}, \qquad
    L_t=\min_j p_{t,j}, \qquad
    C_t=p_{t,n_t}, \qquad
    V_t=\sum_j q_{t,j},
    \label{eq:ohlcv-definition}
\end{equation}

where $n_t$ is the number of transactions in interval $t$, and $O_t$, $H_t$, $L_t$, $C_t$, and $V_t$ are respectively the open, high, low, close, and volume. Their analytical roles are summarized in Table~\ref{tab:ohlcv-meaning}.

\begin{table}[H]
    \centering
    \caption{Meaning of OHLCV observations}
    \label{tab:ohlcv-meaning}
    \begin{tabularx}{\textwidth}{@{}l l X@{}}
        \toprule
        Field & Definition & Information provided \\
        \midrule
        Open & First traded price & Starting valuation of the interval and possible gap from the preceding close. \\
        High & Highest traded price & Upper extreme reached by buyers and sellers during the interval. \\
        Low & Lowest traded price & Lower extreme reached during the interval. \\
        Close & Last traded price & Ending consensus for the interval and the primary input to most indicators. \\
        Volume & Total traded quantity & Level of market participation; it can confirm that a price move occurred with relatively high or low activity. \\
        \bottomrule
    \end{tabularx}
\end{table}

The meaning of a bar depends on its interval. A one-minute bar describes intraday activity, whereas hourly, daily, and weekly bars progressively suppress short-lived fluctuations and emphasize longer trends. Stockrium uses hourly, daily, or weekly observations for short-, medium-, or long-term automated analysis, while its price chart also supports intraday intervals for exploration.

\paragraph{Candlestick and volume charts.}

A candlestick is the conventional visual representation of one OHLC bar. The rectangle, called the \textit{body}, spans the open and close. Its color indicates whether the close is above or below the open. The upper and lower \textit{wicks} extend to the high and low. A long body indicates a large open-to-close movement, whereas a long wick shows that trading extended well beyond the open--close range. Because a candle contains only OHLC values, it does not reveal the order of intrabar movements. The complete high--low range provides a compact view of within-bar dispersion. Volume bars are normally aligned below the price chart so that a breakout, reversal, or quiet consolidation can be interpreted together with market participation. These shapes provide context rather than deterministic rules: the same candle can have different implications depending on the preceding trend, nearby price levels, and volume.

Stockrium displays a candlestick chart by default and can alternatively show a close-price area chart. The chart allows the user to inspect raw market behavior, while overlays and separate panels display derived indicators. Indicators commonly fall into four groups: trend measures such as moving averages, momentum oscillators such as RSI and KDJ, trend--momentum measures such as MACD, and volatility envelopes such as Bollinger Bands. Volume and its moving averages provide supplementary participation context.

\paragraph{Simple and exponential moving averages.}

For an $n$-bar window, the simple moving average (SMA) of closing prices is

\begin{equation}
    \operatorname{SMA}_{n,t}
      =\frac{1}{n}\sum_{i=0}^{n-1} C_{t-i}.
    \label{eq:technical-sma}
\end{equation}

An SMA smooths short-term variation so that the direction of the underlying series is easier to observe. A shorter average responds faster, while a longer average changes more slowly and provides broader trend context. Stockrium uses SMA20 and SMA50. A short average crossing above a long average, a widening positive distance between the two, or the current price remaining above both can support an upward-trend interpretation. Because both averages are calculated from past prices, their signals necessarily lag the market.

The exponential moving average (EMA) places more weight on recent observations:

\begin{equation}
    \operatorname{EMA}_{n,t}
      =\alpha C_t+(1-\alpha)\operatorname{EMA}_{n,t-1},
    \qquad
    \alpha=\frac{2}{n+1}.
    \label{eq:technical-ema}
\end{equation}

In the application-analysis service, the first valid EMA is initialized with the SMA of the first $n$ observations. EMA12, EMA26, and an EMA9 of their difference are used internally by MACD rather than displayed as independent scoring perspectives.

The corresponding volume moving average is

\begin{equation}
    \operatorname{VMA}_{n,t}
      =\frac{1}{n}\sum_{i=0}^{n-1} V_{t-i}.
    \label{eq:technical-vma}
\end{equation}

Stockrium calculates VMA20 and VMA50 so current activity can be compared with its recent baseline. They are displayed with volume and supplied as contextual evidence to the model-backed technical analysis, but they do not add a point to the five deterministic indicator reference signals. Volume is traded quantity and should not be interpreted directly as money flowing into or out of a security.

\paragraph{Relative Strength Index.}

The Relative Strength Index (RSI) is a bounded momentum oscillator. For each close-to-close change, let

\begin{equation}
    \Delta C_t=C_t-C_{t-1}, \qquad
    G_t=\max(\Delta C_t,0), \qquad
    D_t=\max(-\Delta C_t,0),
\end{equation}

where $G_t$ and $D_t$ are non-negative gains and losses. Stockrium uses the 14-bar Wilder smoothing convention. After the initial arithmetic means of 14 changes, the averages are updated as

\begin{align}
    \overline{G}_t
      &=\frac{13\overline{G}_{t-1}+G_t}{14}, \\
    \overline{D}_t
      &=\frac{13\overline{D}_{t-1}+D_t}{14}, \\
    RS_t&=\frac{\overline{G}_t}{\overline{D}_t}, \qquad
    RSI_t=100-\frac{100}{1+RS_t}.
    \label{eq:technical-rsi}
\end{align}

RSI describes the balance between recent upward and downward movement on a scale from 0 to 100; the service returns 100 when the smoothed loss is zero. Values above 70 and below 30 are commonly described as overbought and oversold, but these labels do not themselves predict an immediate reversal; RSI can remain extreme during a strong trend. Stockrium looks for upward momentum through a crossing above 50, an increase within 50--70, or a recovery while below 35.

\paragraph{Moving Average Convergence Divergence.}

Moving Average Convergence Divergence (MACD) compares a faster and a slower EMA:

\begin{align}
    \operatorname{MACD}_t
      &=\operatorname{EMA}_{12,t}-\operatorname{EMA}_{26,t}, \\
    \operatorname{Signal}_t
      &=\operatorname{EMA}_{9}(\operatorname{MACD}_t), \\
    \operatorname{Histogram}_t
      &=\operatorname{MACD}_t-\operatorname{Signal}_t.
    \label{eq:technical-macd}
\end{align}

The MACD line shows whether shorter-term price momentum is stronger or weaker than its longer-term baseline. The signal line smooths that difference, and the histogram makes the distance and its change easier to inspect. Stockrium treats an upward MACD--signal crossing, an increasing histogram while MACD is above the signal line, or a change of the histogram from negative to non-negative as positive momentum evidence. MACD remains a lagging transformation of price and can create repeated false crossings in a sideways market.

\paragraph{Bollinger Bands.}

Bollinger Bands place a volatility-dependent envelope around a moving average. Stockrium uses a 20-bar middle band and two population standard deviations:

\begin{align}
    M_t&=\operatorname{SMA}_{20,t}, \\
    \sigma_t
      &=\sqrt{\frac{1}{20}\sum_{i=0}^{19}(C_{t-i}-M_t)^2}, \\
    U_t&=M_t+2\sigma_t, \qquad
    L_t=M_t-2\sigma_t.
    \label{eq:technical-bollinger}
\end{align}

The band width expands when recent prices are more dispersed and contracts when volatility decreases. The position of price relative to the middle, upper, and lower bands provides trend and overextension context; touching a band is not automatically a buy or sell instruction. Stockrium looks for recovery through the middle band, strengthening movement above the middle band, or an upward reaction near or below the lower band. A price above the upper band is treated cautiously because it may be overextended, although in a strong trend it can also reflect persistent momentum.

\paragraph{KDJ stochastic oscillator.}

KDJ locates the close within the recent high--low range and adds a more responsive $J$ line. For a nine-bar window,

\begin{align}
    HH_{9,t}&=\max_{0\leq i<9} H_{t-i}, \qquad
    LL_{9,t}=\min_{0\leq i<9} L_{t-i}, \\
    RSV_t&=100\frac{C_t-LL_{9,t}}{HH_{9,t}-LL_{9,t}}, \\
    K_t&=\frac{1}{3}RSV_t+\frac{2}{3}K_{t-1}, \\
    D_t&=\frac{1}{3}K_t+\frac{2}{3}D_{t-1}, \qquad
    J_t=3K_t-2D_t.
    \label{eq:technical-kdj}
\end{align}

When the nine-bar range is zero, the application-analysis service sets $RSV_t$ to 50; the first $K$ and $D$ values are also initialized at 50. The $K$ and $D$ lines smooth the location of the close, whereas $J$ magnifies their difference and can move outside the interval $[0,100]$. Stockrium looks for $K$ crossing above $D$, strengthening $J$ while $K>D$, or an upward recovery of $K$ below 20. This sensitivity is useful for short-term turning-point context but also makes KDJ susceptible to noise.

\subsection{Fundamental Analysis}
\label{subsec:fundamental-analysis-foundations}

\paragraph{Concept and purpose.}

Fundamental analysis examines the economic condition of a company and whether its operating performance, financial position, and valuation support its market price and longer-term prospects \cite{penman2013financial}. Its central concepts are \textit{business performance}, including revenue, costs, profit, and cash generation; \textit{financial position}, including assets, liabilities, liquidity, and leverage; \textit{operating efficiency} and \textit{profitability}; \textit{growth and sustainability} across reporting periods; and \textit{valuation}, which relates market price to earnings, book value, cash flow, or enterprise value. These concepts are evaluated through financial statements, ratios, multi-year trends, and comparisons with appropriate peers.

The purpose of fundamental analysis is to assess business quality, financial resilience, earnings capacity, longer-term risk, and whether the market valuation appears reasonable relative to the company's economic performance. It can reveal, for example, that apparent growth is accompanied by deteriorating cash flow, that a high return on equity is driven by leverage, or that a low valuation reflects weak profitability. It therefore supports security selection and medium- to long-term investment judgment. Because accounting data are periodic and backward-looking, fundamental analysis is less suited to precise entry timing and may not immediately capture a newly disclosed event or a rapid change in market expectations.

Its principal inputs are the balance sheet, income statement, and cash-flow statement. The balance sheet reports assets, liabilities, and shareholders' equity at a point in time; the income statement explains how revenue and expenses produce profit over a period; and the cash-flow statement separates operating, investing, and financing cash flows. Because absolute values are affected by company size, analysts commonly convert statement items into ratios and compare them across time and with similar firms. Unless otherwise stated, the ratio equations below are conventional financial-statement definitions summarized from this source rather than formulas created by Stockrium.

\paragraph{Non-financial companies.}

For a non-financial company, Stockrium organizes evidence into liquidity, leverage, operating efficiency, profitability, and valuation. Each category answers a different question and must be interpreted in relation to the company's business model and industry.

Liquidity measures the capacity to meet near-term obligations:

\begin{align}
    \text{Current Ratio}
      &=\frac{\text{Current Assets}}{\text{Current Liabilities}}, \\
    \text{Quick Ratio}
      &=\frac{\text{Current Assets}-\text{Inventories}}
              {\text{Current Liabilities}}, \\
    \text{Cash Ratio}
      &=\frac{\text{Cash and Cash Equivalents}}
              {\text{Current Liabilities}}.
    \label{eq:fundamental-liquidity}
\end{align}

A larger liquidity buffer can reduce short-term payment risk, but an unusually high value can also indicate idle cash, excessive inventory, or inefficient working-capital use. One universal threshold is therefore inappropriate.

Leverage and debt-servicing ratios describe dependence on creditors and the capacity of operating earnings to cover financing costs:

\begin{align}
    \text{Debt-to-Equity}
      &=\frac{\text{Interest-bearing Debt}}
              {\text{Shareholders' Equity}}, \\
    \text{Financial Leverage}
      &=\frac{\text{Total Assets}}
              {\text{Shareholders' Equity}}, \\
    \text{Interest Coverage}
      &=\frac{\text{EBIT}}{\text{Interest Expense}}.
    \label{eq:fundamental-leverage}
\end{align}

Leverage can increase shareholder returns when operations perform well, but it also magnifies losses and refinancing risk. Higher interest coverage generally provides a larger buffer for interest payments. Normal levels differ greatly between industries, so leverage should be compared with an appropriate peer group.

Operating-efficiency ratios examine how effectively assets and working capital generate revenue:

\begin{align}
    \text{Asset Turnover}
      &=\frac{\text{Revenue}}{\text{Average Total Assets}}, \\
    \text{Fixed-Asset Turnover}
      &=\frac{\text{Revenue}}{\text{Average Net Fixed Assets}}, \\
    \text{Days Inventory}
      &=\frac{\text{Average Inventory}}{\text{Cost of Goods Sold}}\times365, \\
    \text{Days Receivable}
      &=\frac{\text{Average Accounts Receivable}}{\text{Net Revenue}}\times365.
    \label{eq:fundamental-efficiency}
\end{align}

Higher turnover generally means that each unit of assets supports more revenue. Increasing inventory or receivable days may indicate slow-moving stock or slower collection, although the expected level depends on the operating cycle. A manufacturer, utility, retailer, and software company should not be expected to have identical ratios.

Profitability ratios relate earnings to sales, assets, or shareholder capital:

\begin{align}
    \text{Gross Margin}
      &=\frac{\text{Gross Profit}}{\text{Revenue}}, \\
    \text{Net Margin}
      &=\frac{\text{Net Income}}{\text{Revenue}}, \\
    ROA&=\frac{\text{Net Income}}{\text{Average Total Assets}}, \\
    ROE&=\frac{\text{Net Income}}
              {\text{Average Shareholders' Equity}}.
    \label{eq:fundamental-profitability}
\end{align}

Gross margin reflects pricing and direct production or service-delivery cost, while net margin also incorporates operating expenses, financing costs, and taxes. ROA measures profit relative to the asset base and ROE measures profit relative to shareholder capital. A high ROE is not necessarily created by strong operations because it can be amplified by a small, highly leveraged equity base. The three-part DuPont identity separates these sources:

\begin{equation}
    ROE=
    \underbrace{\frac{\text{Net Income}}{\text{Revenue}}}_{\text{Net Margin}}
    \times
    \underbrace{\frac{\text{Revenue}}{\text{Average Total Assets}}}_{\text{Asset Turnover}}
    \times
    \underbrace{\frac{\text{Average Total Assets}}
                     {\text{Average Shareholders' Equity}}}_{\text{Financial Leverage}}.
    \label{eq:fundamental-dupont}
\end{equation}

This standard DuPont decomposition \cite{openstaxDupont2022} reveals whether returns are supported mainly by margin, efficient asset use, or leverage. The direction of each component over several years is more informative than one isolated ROE value.

Valuation ratios compare the market value assigned to the company with earnings, book value, or operating cash-generating capacity:

\begin{align}
    EPS
      &=\frac{\text{Earnings Attributable to Ordinary Shareholders}}
              {\text{Weighted-average Ordinary Shares}}, \\
    P/E&=\frac{\text{Market Price per Share}}{EPS}, \\
    P/B&=\frac{\text{Market Price per Share}}{\text{Book Value per Share}}, \\
    EV/EBITDA&=\frac{\text{Enterprise Value}}{EBITDA}.
    \label{eq:fundamental-valuation}
\end{align}

A lower multiple can indicate a cheaper relative valuation, but it can also reflect low expected growth or high risk. P/E is not meaningful with negative earnings and EV/EBITDA is difficult to interpret with non-positive EBITDA. Stockrium therefore treats valuation as comparative evidence rather than proof that a security is under- or overvalued.

\paragraph{Financial-sector companies.}

Bank analysis requires a different balance-sheet interpretation because deposits and loans are core operating inputs and outputs rather than ordinary financing and receivable items. Stockrium identifies an Industry Classification Benchmark 8xxx company as financial and organizes its available evidence using CAMELS-related categories. Because this group also contains insurers, securities firms, and other non-bank companies, specialist banking fields may be unavailable and general ratios are then used as proxies; the result is not a complete supervisory CAMELS rating.

\begin{description}
    \item[\textbf{Capital adequacy.}]
    Capital provides a buffer against unexpected losses. For a bank, the principal measure is
    \[
        CAR=\frac{\text{Regulatory Capital}}{\text{Risk-weighted Assets}}.
    \]
    Stockrium also considers financial leverage, debt-to-equity, equity-to-assets, equity-to-liabilities, equity-to-loans, and book value per share when available.

    \item[\textbf{Asset quality.}]
    Credit quality and its reserve protection can be summarized by
    \[
        NPL=\frac{\text{Non-performing Loans}}{\text{Gross Loans}},
        \qquad
        \text{NPL Coverage}
          =\frac{\text{Loan-loss Reserves}}{\text{Non-performing Loans}}.
    \]
    A lower NPL ratio generally indicates better loan quality, while higher coverage provides a stronger reserve buffer. Loans-to-equity, ROA, reserves-to-loans, provision expense, and credit growth add context.

    \item[\textbf{Management efficiency.}]
    The cost-to-income ratio is
    \[
        CIR=\frac{\text{Operating Expenses}}{\text{Operating Income}}.
    \]
    A lower CIR generally reflects better cost control. Stockrium supplements it with asset turnover, ROIC, and non-interest income relative to interest income.

    \item[\textbf{Earnings.}]
    Net interest margin is conventionally
    \[
        NIM=\frac{\text{Net Interest Income}}{\text{Average Earning Assets}},
    \]
    where net interest income equals interest income minus interest expense. ROE, ROA, net margin, EPS, yield on earning assets, cost of funds, and the current-and-savings-account (CASA) deposit ratio provide additional profitability and funding context.

    \item[\textbf{Liquidity.}]
    The loan-to-deposit ratio is
    \[
        LDR=\frac{\text{Gross Customer Loans}}{\text{Customer Deposits}}.
    \]
    A very high value may indicate a narrow funding-liquidity buffer, whereas a very low value can indicate underused funding. Deposit growth and available cash, quick, and current ratios provide supplementary evidence.

    \item[\textbf{Sensitivity to market risk.}]
    This category concerns exposure to interest-rate, exchange-rate, and other market changes. The current dataset has no dedicated quantitative sensitivity field, so Stockrium reports the limitation rather than inferring a measure.
\end{description}

For financial-sector valuation, Stockrium uses P/B because book equity is particularly relevant to balance-sheet-driven businesses and excludes P/E and EV/EBITDA from this analytical output. P/B must nevertheless be read with profitability, asset quality, growth, and comparable institutions; a low P/B alone does not establish undervaluation.

\paragraph{Trend and industry context.}

One annual observation is only a snapshot. Stockrium retains up to the five latest annual observations and orders them chronologically. When both endpoints are positive, it summarizes their change using the standard compound annual growth rate definition \cite{damodaranHistoricalGrowth2026}:

\begin{equation}
    \operatorname{CAGR}(x)
      =\left(\frac{x_n}{x_0}\right)^{\frac{1}{y_n-y_0}}-1,
    \label{eq:fundamental-cagr-foundation}
\end{equation}

where $x_0$ and $x_n$ are the earliest and latest valid observations and $y_n-y_0$ is the number of elapsed years. CAGR is omitted when fewer than two observations exist, the period is non-positive, or either endpoint is non-positive. It smooths the path between two endpoints and should be read alongside the annual sequence. For measures expressed in days, a negative CAGR can be favorable because inventory or collection periods have shortened.

The system also compares a company with annual medians for its ICB industry. A median is less sensitive than a mean to extreme values and is shown only when the industry data contain at least three peers. For non-financial firms, the comparison includes profitability, margins, asset turnover, leverage, P/E, and P/B. For financial-sector firms, it includes available profitability, leverage, liquidity, P/B, CAR, NIM, LDR, and NPL measures. The comparison gives a ratio meaning only in relation to the economics and normal capital structure of similar firms.

\subsection{Financial News Analysis}
\label{subsec:financial-news-analysis-foundations}

\paragraph{Concept and purpose.}

Financial news analysis examines time-stamped qualitative evidence about events and changing expectations that may affect a security before their full consequences appear in financial statements. Its principal concepts are \textit{relevance}, whether an item materially concerns the analyzed company, industry, or market; \textit{event and transmission mechanism}, what happened and how it may alter revenue, costs, financing, risk, or valuation; \textit{sentiment}, the directional implication for a specified entity rather than the emotional tone of the wording; \textit{timeliness and novelty}, whether the information was available and not already stale or duplicated at the decision time; and \textit{provenance and credibility}, which distinguish official disclosures and verified reporting from opinion or unverified claims.

The purpose of financial news analysis is to detect recent catalysts and risks, update an assessment between formal reporting periods, and explain why market expectations or price behavior may be changing. It can identify information that historical indicators and annual statements do not yet contain, such as a new contract, policy change, management event, litigation, or earnings surprise. It also supports event-driven monitoring and helps determine whether existing technical and fundamental conclusions remain valid. News alone is insufficient because an item may be immaterial, misinterpreted, already reflected in price, or relevant only over a different horizon; it therefore requires source, temporal, and company-specific context.

Company-level news includes earnings announcements, investment and financing decisions, management changes, contracts, litigation, and regulatory disclosures. Industry and macroeconomic news includes sector demand, interest rates, exchange rates, inflation, and government policy. Relevance is more specific than the presence of a company name or symbol: the analyst must identify the affected entity, the event, the mechanism through which it may alter revenue, costs, financing, risk, or valuation, and the horizon over which that effect may occur. Source provenance should also be retained because an official disclosure, a reported fact, an analyst opinion, and an unverified claim do not provide equivalent evidence.

\paragraph{Summarization.}

News summarization reduces a long article to the information needed for analysis. A useful financial summary should preserve the affected company or industry, the principal event, material figures and dates, the expected transmission mechanism, and important uncertainty. It should distinguish completed events from forecasts or proposals and facts from the author's interpretation. Abstractive summaries are more readable than copied excerpts, but they can omit qualifications or introduce unsupported relationships. A generated summary should therefore remain traceable to the original article rather than replace it as the authoritative source.

\paragraph{Sentiment analysis.}

Financial sentiment analysis assigns a directional interpretation such as positive, neutral, or negative. In this setting, sentiment represents the expected implication for a specified company or market, not merely the emotional tone of the words. The same event can benefit one company, harm a competitor, and have an uncertain sector-wide effect; apparently positive language may also describe a result below market expectations. Domain-specific work such as FinBERT and Vietnamese stock-title classification demonstrates why financial vocabulary and local-language context matter \cite{araci2019finbert,tung2021phobert}. Mixed, conditional, or immaterial information should remain neutral instead of being forced into a directional trading signal.

\subsection{Rationale for Combining Technical, Fundamental, and News Analysis}
\label{subsec:combined-financial-analysis-foundations}

The three forms of analysis answer related but different decision questions. Technical analysis asks \textit{what the market is doing and when conditions may be timely}; fundamental analysis asks \textit{what economic performance and valuation support the security over a longer horizon}; and news analysis asks \textit{what has recently changed and how that change may affect expectations}. Their data frequencies also differ: market prices and volume can change continuously, news arrives irregularly when events occur, and financial statements are published periodically. Combining them gives the decision process evidence about market behavior, underlying business condition, and new information instead of treating any one perspective as a complete representation of the security.

The combination is necessary because the principal limitation of one perspective is often addressed by another. Technical indicators can reveal momentum but cannot determine whether it is supported by business performance or caused by a temporary event. Fundamental analysis can identify a financially strong and reasonably valued company but cannot establish that the current market trend is favorable or that an adverse event has just occurred. News can identify a catalyst or risk quickly but cannot by itself show whether the effect is financially material, already reflected in the price, or persistent. Agreement among the three perspectives provides stronger cross-domain support, while disagreement is itself useful evidence: strong fundamentals with weak technical conditions may support waiting for better timing, and positive momentum accompanied by adverse material news may indicate elevated reversal risk.

Stockrium operationalizes this complementarity by assigning technical, fundamental, and news evidence to specialized components and synthesizing their outputs according to the selected investment horizon. Shorter horizons emphasize current market behavior and recent catalysts, whereas longer horizons give greater importance to business performance and valuation. The final result remains conditional because the three inputs are neither complete nor statistically independent: news can affect price, price can change valuation ratios, and all sources can contain delay or measurement error. Their combination reduces analytical blind spots and improves the traceability of a recommendation, but it does not eliminate uncertainty or guarantee investment returns.

The relative importance assigned to these sources is therefore an explicit analytical assumption rather than a universally optimal formula. Stockrium exposes automatic and supported manual configurations so that their inputs and resulting component scores can be inspected. Stockrium Lab extends this transparency by retaining experiment configurations and outputs for side-by-side comparison. Such comparison can reveal sensitivity, disagreement, or unexpected behavior that warrants further investigation; it does not by itself identify a configuration that will remain profitable in future market conditions.
